\chapter{Upper bounds}\label{chap:upper-bounds}

Before we proceed towards seeing techniques for lower bounds, it would be good to be aware of some surprising upper bounds for polynomials. 
It is often the case that many lower bound proofs are really about leveraging some interesting upper bounds, and we'll see several examples of this in the survey. 


\section{Determinant}

In \autoref{thm:det-abp}, we saw that the $n\times n$ symbolic determinant can be computed by an algebraic branching program of size $\poly(n)$. 
We record that result again here as that is certainly an example of a non-trivial upper bound. 

\begin{theorem*}[(restatement of \autoref{thm:det-abp})] The $n\times n$ symbolic determinant can be computed by an algebraic branching program of depth $n$ and width $O(n^3)$. 
\end{theorem*}

\section{Permanent}

Although the $n\times n$ symbolic permanent $\Perm_n$ is believed to be hard to compute, there are non-trivially small computations of the permanent. 

The following computation is just `inclusion-exclusion' to compute $\Perm_n$, and is called Ryser's formula. 

\begin{lemma}[Ryser's formula] For any $n \in \N$, 
  \[
  \operatorname{perm}(X) = \Perm_n = \sum_{S \subseteq [n]} (-1)^{n-|S|} \prod_{i=1}^n \sum_{j\in S} x_{ij}.
  \]
  Thus, $\Perm_n$ can be computed by a \emph{homogeneous, set-multilinear} $\Sigma\Pi\Sigma$ circuit of size $2^n \cdot \poly(n)$. 
\end{lemma}

If the symbolic matrix $X$ on the LHS is replaced by the matrix each of whose rows are $(x_1, \ldots, x_n)$, then we obtain the following identity to compute a monomial

\begin{lemma}[Fischer's identity]
\[
n! \cdot x_1 \cdots x_n = \sum_{S \subseteq [n]} (-1)^{n - |S|} \inparen{\sum_{j \in S} x_j}^n.
\]
\end{lemma}

In fact, one can also construct an ABP for $\Perm_n$ (with edges carrying linear forms) consisting of just $2^n$ vertices. 

\begin{exercise}
Show that there is an ABP for $\Perm_n$ consisting of $2^n$ vertices. 

(Hint: Each node in the ABP will be indexed by a subset of $[n]$, with edges going from nodes corresponding to sets of size $i$ to those corresponding to sets of size $i+1$.)
\end{exercise}

\section{Elementary symmetric polynomials}

The \emph{elementary symmetric polynomial} $\ESym_d(x_1,\ldots, x_n)$ is defined as 
\[
\ESym_d(x_1,\ldots, x_n) = \sum_{\substack{S \subset [n]\\|S| = d}} \prod_{i\in S} x_i,
\]
which is the sum of all multilinear monomials corresponding to size $d$ subsets of $[n]$. 

\subsection{Some straightforward computations}

The standard way of computing this as a circuit is just a simple recursive method via 
\[
\ESym_{d}(x_1,\ldots, x_n) = \ESym_{d}(x_1,\ldots, x_{n-1}) + x_n \ESym_{d-1}(x_1,\ldots, x_{n-1}).
\]
As a formula, this has size $\binom{n}{d}$ but one can show (with some care) that this yields a circuit of size $\poly(n,d)$. 


\noindent
Another computation is 
\[
\ESym_{d}(x_1,\ldots, x_n) = \sum_{j=0}^d \ESym_{j}(x_1,\ldots, x_{n/2}) \cdot \ESym_{d-j}(x_{n/2 + 1},\ldots, x_{n}).
\]
With a bit of care, one can prove the following. 

\begin{lemma}
The above expression for the elementary symmetric polynomial yields an algebraic formula of size $n \cdot \binom{d + \log n}{d} \leq n \cdot 2^{d + \log n} = n^2 \cdot 2^d$. 
\end{lemma}

Both of these computations are homogeneous computations but of large depth. 

\subsection{Depth-3 circuits}

This polynomial \emph{feels like} an algebraic analogue of the boolean ``Threshold'' gate (that returns $1$ if at least $d$ of its input bits are $1$) and one might think that this is perhaps not computable by $\poly(n,d)$-size constant-depth circuits. An important technique we will see in the next chapter is that of \emph{interpolation} (\autoref{lem:interpolation}). 

\begin{lemma}[Interpolation (informal)] For any set of distinct field elements $\alpha_0, \ldots, \alpha_d$ and integers $0 \leq i \leq d$, there are field elements $\beta_{i,0}, \cdots, \beta_{i,d}$ such that for all polynomial $f(x) = f_0 + f_1 x + \cdots + f_d x^d$ we have
  \[
  f_i = \sum_{j=0}^d \beta_{i,j} f(\alpha_j).
  \]
  That is, we can ``interpolate'' any coefficient from sufficiently many evaluations of $f$. 
\end{lemma}

With this, consider the polynomial $f(x_1,\ldots, x_n, t) = (1 + tx_1) \cdots (1 + tx_n)$. Interpreted as a univariate in $t$ (with coefficients as polynomials in $x_1,\ldots, x_n$), the coefficient of $t^i$ is precisely $\ESym_i$. Therefore, we immediately get the following. 

\begin{lemma}[Depth-$3$ circuit for $\ESym_d$ over any large field]\label{lem:esym-benor-trick}
  For all $d \leq n$, and any field $\F$ with $|\F| \geq n+1$, there is a circuit for $\ESym_d$ of the form
  \[
  \ESym_d(x_1,\ldots, x_n) = \sum_{j=0}^d  \beta_{j} \cdot (1 + \alpha_j x_1) \cdots (1 + \alpha_j x_n).
  \]
\end{lemma}
The above is a non-homogeneous circuit of depth $3$ for the elementary symmetric polynomial. \\

There are several other ways of computing the elementary symmetric polynomial 

\subsection{Read-once oblivious branching programs}

One other way to construct the elementary symmetric polynomial is via the following branching program:
\[
\ESym_d(x_1,\ldots, x_n) = \inparen{\begin{bmatrix} 
  1 & x_1 & 0   & \cdots & 0\\
    & 1   & x_1 & \cdots & 0\\
    &    &  1 & \ddots & \vdots\\
    & &  & \ddots & \vdots\\
  & &  &  & 1
\end{bmatrix}\cdot \begin{bmatrix} 
  1 & x_2 & 0   & \cdots & 0\\
    & 1   & x_2 & \cdots & 0\\
    &    &  1 & \ddots & \vdots\\
    & &  & \ddots & \vdots\\
  & &  &  & 1
\end{bmatrix}
 \cdots \begin{bmatrix} 
  1 & x_n & 0   & \cdots & 0\\
    & 1   & x_n & \cdots & 0\\
    &    &  1 & \ddots & \vdots\\
    & &  & \ddots & \vdots\\
  & &  &  & 1
\end{bmatrix}}_{(1,d+1)}
\]
That is, we can write $\ESym_d(x_1,\dots, x_n)$ as a single entry of a product of $(d+1)\times(d+1)$ matrices, each matrix being a \emph{univariate} in a single variable. Such ABPs are also called \emph{read-once oblivious ABPs}, and we will encounter them later in the survey.

\subsection{Newton-Girard identities}

The Newton-Girard identities, over characteristic zero fields, allow one to write the $\ESym$ polynomials in terms of the $\PSym$ polynomials
\[
\PSym_d(x_1,\ldots, x_n) = x_1^d + \cdots + x_n^d.
\]
\begin{lemma}[Newton-Girard identities]
For all $0 \leq d \leq n$, we have
\[
d \cdot \ESym_d + \sum_{i=1}^d (-1)^{d-i}\PSym_i \ESym_{d-i} = 0
\]
\end{lemma}
\begin{proof}
  We'll first prove the statement for $d = n$. 
  Let $f(t) = (t - x_1) \cdots (t - x_n) = \sum t^i (-1)^{n-i} \ESym_{n-i}$. Clearly $f(x_1) = f(x_2) = \cdots = f(x_n) = 0$ and hence $\sum_i f(x_i) = 0$. 
  \begin{align*}
  0 = \sum_i f(x_i) &= \sum_{i=0}^d (-1)^{n-i} \PSym_i \cdot \ESym_{n-i}\\
    & = n \cdot \ESym_n + \sum_{i=1}^n (-1)^{n-i} \PSym_i \cdot \ESym_{n-i}
  \end{align*}
  (where $\ESym_0 = 1$ by convention for all $n$, and $\PSym_0$ is usually not defined so that's handled separately). 

  For $d < n$, consider the LHS and RHS of the purported equality
  \[
  d \cdot \ESym_d = \sum_{i=1}^d (-1)^{d-i} \PSym_i \ESym_{d-i}
  \]
  For every monomial on the LHS, we can set the other $(n-d)$ variables to zero and notice that the RHS is the Newton identities for $n' = d$ and hence holds. Since the coefficients match for all monomials of degree $d$, the above must indeed be an equality.
\end{proof}

The Newton-Girard identities are also expressible as a matrix
\[
  \ESym_d = \frac{1}{d!} \cdot 
  \begin{bmatrix}\PSym_1 & 1 & 0 & \cdots & 0 & 0\\ 
    \PSym_2 & \PSym_1 & 2 & \cdots & 0 & 0 \\ 
    \vdots& \vdots & \ddots & \ddots &  & \vdots \\
    \vdots& \vdots &  & \ddots & \ddots & \vdots \\
    \PSym_{d-1} & \PSym_{d-2} & \PSym_{d-3}& \cdots & \PSym_1 & d-1 \\ 
    \PSym_d & \PSym_{d-1} & \PSym_{d-2} & \cdots & \PSym_2 & \PSym_1 
  \end{bmatrix}.
\]
The RHS, when expanded into monomials yields a homogeneous (why?) expression of the form 
\[
\ESym_d = \sum_{\substack{0\leq e_1,\ldots, e_d\leq d\\ e_1 + 2e_2 + \cdots + de_d = d}} c_{\vece} \cdot \PSym_1^{e_1} \cdot \PSym_2^{e_2} \cdots \PSym_d^{e_d}
\]
The number of summands equals the number of partitions of $d$, and turns out this is bounded by $2^{O(\sqrt{d})}$. This then yields the following. 

\begin{lemma}[Homogeneous depth-4 circuit for $\ESym_d$]\label{lem:hom-d4-ckt-esym}
Over any field of zero (or large) characteristic, there is a homogeneous depth-$4$ circuit of size $2^{O(\sqrt{d})} \cdot \poly(n)$ for $\ESym_d(x_1,\ldots, x_n)$. 
\end{lemma}

\begin{exercise}
Use the Newton identities to get a circuit for $\Det_n$ by using the fact that $\Det_n = \lambda_1 \cdots \lambda_n$ where $\lambda_i$'s are the eigenvalues of the symbolic matrix $X$. 
\end{exercise}

\section{Rectangular permanents}

Much of this section is from Forbes~\cite{TODO} that will be discussed in detail in TODO. 

Consider the following polynomial called the ``rectangular permanent'', denoted by $\RPerm_{r\times n}$, for positive integers $r \leq n$. In words, $\RPerm_{r\times n}$ is the polynomial obtained by considering an $r\times n$ matrix of distinct indeterminates, and adding all monomials that can be obtained by picking $r$ variables from distinct rows and columns. More formally, 
\[
\RPerm_{r,n}(X_{r\times n}) = \sum_{\sigma: [r] \hookrightarrow [n]} \prod_{i=1}^r x_{i, \sigma(i)}
\]
where $\sigma:[r] \hookrightarrow [n]$ denotes an injective function. 

An obvious computation is of course to just write it as 
\[
\sum_{S \in \binom{[n]}{r}}\Perm_n(X_S)
\]
which has size $n^{O(r)} \cdot 2^{O(r)}$. Can we do better, particularly in regimes when $r \ll n$? For example, when $r = 2$, we have
\begin{align*}
\RPerm_{2,n} & := \sum_{1 \leq i,j\leq n, i\neq j} x_{1i} x_{2_j}\\
            &  = \inparen{\sum_{1 \leq i \leq n} x_{1i}} \cdot \inparen{\sum_{1 \leq j \leq n} x_{2j}} - \inparen{\sum_{1 \leq i \leq n} x_{1i}x_{2i}}
\end{align*}
which only has size $O(n)$. This can be generalised for larger $r$ and is called the Binet-Minc identity. 

\begin{lemma}[Circuits for rectangular permanents]\label{lem:rectangular-perm-computation}
  For all $1 \leq r \leq n$, we can write $\RPerm_{r,n}$ as
  \[
  \RPerm_{r,n}(X_{r\times n}) = \sum_{\lambda \in \mathcal{P}([r])} (-1)^{r - \abs{\lambda}} \prod_{S \in \lambda} (|S| - 1)! \cdot \inparen{\sum_{j=1}^n \sum_{i\in S} x_{ij}}. 
  \]
  In particular, the above is a depth-$4$ circuit of size $B_r \cdot \poly(r) \cdot n$ where $B_r$ is the $r$-th Bell number ($B_r = O((r / \log r)^r)$) that counts the number of distinct partitions of $[r]$. 
\end{lemma}

There are several proofs of this, but we'll give a slightly not-so-self-contained proof via the M\"obius inversion formula that may be applicable in other contexts. 

\paragraph{General context of M\"obius inversion:} Consider any partially ordered set $(\mathcal{P}, \leq)$ (for example the set of subsets of $[n]$ under inclusion, or the set of partitions of $[n]$ under refinement, etc.). Given such a set, we can define a unique \Mobius{} function $\mu: \mathcal{P}\times \mathcal{P} \rightarrow \Z$ that satisfies the following properties:
\begin{align*}
\text{For all $A \in \mathcal{P}$,}& &\mu(A, A)  &= 1.\\
\text{For all $A\leq B \in \mathcal{P}$,}& &\sum_{A \leq T \leq B} \mu(A, T) &= 0.
\end{align*}
\noindent
Equivalently, 
\[
\mu(A,B) = \begin{cases}
  1 & \text{if $A = B$}\\
  0 & \text{if $A \nleq B$}\\
  -\sum_{A \leq T < B} \mu(A, T) & \text{otherwise}.
  \end{cases}
\]
(As an example, for the poset of subsets of $[n]$ under inclusion, $\mu(A,B) = (-1)^{|B \setminus A|}$ when $A \subseteq B$.)

The \Mobius{} function is set up to make the following statement hold. 

\begin{lemma}[\Mobius{} inversion formula]
  Suppose $f,g: \mathcal{P} \rightarrow R$ that satisfies
  \[
  f(A) = \sum_{B \geq A} g(B).
  \]
  Then, 
  \[
  g(A) = \sum_{B \geq A} \mu(A,B) \cdot f(B).
  \]
\end{lemma}

With this context, we can now move to rectangular permanents. 

\begin{proof}[Proof of \cref{lem:rectangular-perm-computation}]
  For a function $h:[r]\rightarrow [n]$, let $\ker h$ denote the partition of $[r]$ defined by setting $i \sim j$ if and only if $h(i) = h(j)$. If $h$ is one-to-one, then $\ker h$ is the finest partition $\hat{0}$ consisting of singletons. 

  Consider the poset $\mathcal{P}$ consisting of partitions of $[r]$ where $\lambda \geq \tau$ if $\tau$ is a finer partition of $\lambda$. Let us define $g: \mathcal{P} \rightarrow \F[\vecx]$ given by 
  \begin{align*}
    g(\lambda) & = \sum_{\substack{h: [r] \rightarrow [n]\\ \ker h = \lambda}} \prod_{i=1}^r x_{i, h(i)}\\
    \implies g(\hat{0}) & = \RPerm_{r,n}(X_{r\times n}).
  \end{align*}
  Let us define $f: \mathcal{P} \rightarrow \F[\vecx]$ given by $f(\lambda) = \sum_{\tau \geq \lambda} g(\tau)$.
  \begin{align*}
  f(\lambda) & = \sum_{\substack{h: [r] \rightarrow [n]\\ \ker h \geq \lambda}} \prod_{i=1}^r x_{i, h(i)} = \prod_{S \in \lambda} \sum_{j=1}^m \prod_{i\in S} x_{i, j}
  \end{align*}
  By the \Mobius{} inversion formula, we now have
  \begin{align*}
    \RPerm_{r,n}(X) = g(\hat{0}) & = \sum_{\lambda \geq \hat{0}} \mu(\hat{0}, \lambda) \cdot f(\lambda)\\
                                 & = \sum_{\lambda \geq \hat{0}} \mu(\hat{0}, \lambda) \cdot \prod_{S \in \lambda} \sum_{j=1}^m \prod_{i\in S} x_{i, j}.
  \end{align*}
  Turns out $\mu(\hat{0}, \lambda) = (-1)^{r - |\lambda|} \prod_{S \in \lambda} (|S| - 1)!$. Plugging this into the above expression gives the claim. 
\end{proof}

\begin{exercise}
Formally prove the \Mobius{} inversion formula.

For the partition lattice of $[r]$, show that $\mu(\hat{0}, \lambda) = (-1)^{r - |\lambda|} \prod_{S \in \lambda} (|S| - 1)!$. More generally, what is $\mu(\lambda, \tau)$ for $\lambda \leq \tau$?
\end{exercise}


The above expression for the rectangular permanent also provides a homogeneous depth-$4$ circuit for the standard permanent of size $B_n \cdot \poly(n) = \inparen{\frac{n}{\log n}}^{O(n)}$. Similarly, if we were to work with the $r \times n$ matrix $Y$ with $Y_{ij} = x_j$ (i.e., all rows are identical), then 
\[
\RPerm_{r,n}(Y) = d! \cdot \ESym_d(x_1,\ldots, x_n),
\]
and we get a $\Sigma\Pi\Sigma\wedge$ circuit of size $p(n)$ (the number of partitions of the integer $n$, rather than the number of partitions of the set $[n]$) and hence we get a $2^{O(\sqrt{n})}$-sized depth-$4$ circuit for $\ESym_d$, recovering the bound in \cref{lem:hom-d4-ckt-esym}. 


%%% Local Variables: 
%%% mode: latex
%%% TeX-master: "fancymain"
%%% End: 
